\documentclass[rapids]{jfm}

\usepackage{amsmath,amssymb}
\usepackage{booktabs}

\shorttitle{Approximate conservation laws for point vortex dynamics}
\shortauthor{B. Sanchez}

\title{A broad family of approximate conservation laws\\ for point vortex dynamics}

\author{Bryan Sanchez\corresp{Email address for correspondence: bryansanchez@proton.me}}
\affiliation{\aff{1}Independent researcher}

\begin{document}

\maketitle

\begin{abstract}
For any smooth function~$f$, the pairwise-weighted quantity $Q_f = \sum_{i<j} \Gamma_i \Gamma_j \, f(r_{ij})$ is approximately conserved during Kirchhoff evolution of two-dimensional point vortices, where $\Gamma_i$ are circulations and $r_{ij}$ inter-vortex distances. Numerical integration across 3- to 8-vortex systems, including chaotic configurations, confirms $\sigma_Q/|\bar{Q}| < 10^{-4}$ for well-behaved~$f$. Two members reduce to known exact invariants: $f = \ln r$ gives the Hamiltonian~$H$, and $f = r^2$ gives the angular impulse~$L_z$ (times total circulation). The optimal non-trivial member is $f = \sqrt{r}$, with fractional variance $3 \times 10^{-11}$ in the restricted three-vortex problem. In three dimensions, $Q_{1/r}$ achieves the best conservation; both optimal choices are Green's functions of the Laplacian and equal the kinetic energy up to constants. The family exhibits a dichotomy: members with $f(r) \to \infty$ as $r \to 0$ detect vorticity concentration, while members with $f(r) \to 0$ as $r \to 0$ resist stretching. No single~$f$ achieves both. The triplet generalisation $\sum_{i<j<k} \Gamma_i \Gamma_j \Gamma_k \,\mathrm{Area}(ijk)$ is not conserved, confirming the pairwise structure is special.
\end{abstract}

\begin{keywords}
vortex dynamics, conservation laws, vortex interactions
\end{keywords}

\section{Introduction}

Point vortex dynamics, introduced by \citet{helmholtz1858} and formalised by \citet{kirchhoff1876}, describes the motion of $N$ singular vortices in two dimensions. Each vortex~$i$ carries circulation $\Gamma_i$ and moves according to the velocity field induced by all other vortices. The equations of motion are Hamiltonian, and four classical conservation laws are known: the Hamiltonian~$H$ (kinetic energy), two components of linear impulse $P_x$, $P_y$, and the angular impulse~$L_z$.

For $N \geq 4$ vortices with generic circulations, these four invariants do not suffice for integrability \citep{newton2001, modin2021}, and no other independent pairwise invariants are known. Standard surveys \citep{aref1983, saffman1992, newton2001, aref2007, newton2014} list only these four, with no mention of approximate pairwise families.

I report that the pairwise-weighted quantity
\begin{equation}
Q_f = \sum_{i<j} \Gamma_i \Gamma_j \, f(r_{ij})
\label{eq:qf}
\end{equation}
is approximately conserved for a broad, continuous family of smooth functions~$f$, where $r_{ij} = |z_i - z_j|$ is the distance between vortices $i$ and~$j$. This is not a perturbative result valid only near integrable configurations: it holds for generic circulations, random initial conditions, and chaotic trajectories with $N$ up to~8.

Two members are exactly conserved: $f(r) = \ln r$ gives $Q_f \propto H$, and $f(r) = r^2$ gives $Q_f = \Gamma_\mathrm{tot} L_z - |P|^2$. The remaining members are new. In three dimensions, $Q_{1/r}$ is optimal. Both $-\ln r$ (2D) and $1/r$ (3D) are Green's functions of the Laplacian, and both equal the kinetic energy. I conjecture this holds in any dimension.

\section{Setup}

Consider $N$ point vortices with positions $z_i(t) \in \mathbb{C}$ and circulations $\Gamma_i \in \mathbb{R} \setminus \{0\}$. The equations of motion are
\begin{equation}
\frac{\mathrm{d}\bar{z}_i}{\mathrm{d}t} = \frac{1}{2\pi \mathrm{i}} \sum_{j \neq i} \frac{\Gamma_j}{z_i - z_j},
\end{equation}
with Hamiltonian $H = -(1/2\pi) \sum_{i<j} \Gamma_i \Gamma_j \ln r_{ij}$. Conservation quality is measured by fractional variance $\sigma_Q / |\bar{Q}|$ evaluated over the trajectory; values below $5 \times 10^{-3}$ pass the acceptance threshold. The classical invariants achieve $\sigma_Q / |\bar{Q}| < 10^{-12}$, limited by integrator tolerance.

Integration uses the Dormand--Prince RK4(5) method \citep{dormand1980} with absolute and relative tolerances $10^{-12}$ over $T = 10$ time units (2D) or $T = 2$ (3D filaments). Seven configurations span the range from integrable to chaotic: restricted 3-vortex ($\Gamma_3 = \epsilon$), equal 3-vortex, generic 4-vortex ($\Gamma_i = \{1, 2, -1, 0.5\}$, chaotic), hierarchical 5-, 6-, 7-vortex, and an 8-vortex dipole array.

\section{Results}

\subsection{Universal conservation across $N$}

Table~\ref{tab:universal} shows fractional variance for $Q_f$ across all configurations. Every tested~$f$ achieves fractional variance well below the $5 \times 10^{-3}$ threshold for $N = 3$ through~8.

\begin{table}
\centering
\caption{Fractional variance of $Q_f$ for selected configurations and functions~$f$. Entries marked $< 10^{-20}$ are exact to machine precision.}
\label{tab:universal}
\begin{tabular}{@{}llcccc@{}}
\toprule
$N$ & Configuration & $f = r$ & $f = r^2$ & $f = \sqrt{r}$ & $f = \ln r$ \\
\midrule
3 & Restricted  & $2 \times 10^{-10}$ & $< 10^{-20}$ & $3 \times 10^{-11}$ & $< 10^{-20}$ \\
3 & Equal       & $3 \times 10^{-8}$  & $< 10^{-20}$ & $4 \times 10^{-9}$  & $< 10^{-20}$ \\
4 & Generic     & $2 \times 10^{-6}$  & $< 10^{-20}$ & $3 \times 10^{-7}$  & $< 10^{-20}$ \\
5 & Hierarchical & $2 \times 10^{-7}$ & $< 10^{-20}$ & $3 \times 10^{-8}$  & $< 10^{-20}$ \\
8 & Dipole array & $1 \times 10^{-4}$ & $3 \times 10^{-6}$ & $1 \times 10^{-6}$ & $< 10^{-20}$ \\
\bottomrule
\end{tabular}
\end{table}

\subsection{Optimal power law}

Table~\ref{tab:power} shows fractional variance for $Q_n = \sum \Gamma_i \Gamma_j\, r_{ij}^n$ across a range of exponents in the restricted 3-vortex problem. The optimal non-trivial exponent is $n = 0.5$ (i.e.\ $f = \sqrt{r}$). Conservation quality degrades with $|n|$, but all tested powers from $n = -2$ to $n = 4$ pass the threshold.

\begin{table}
\centering
\caption{Conservation quality versus power-law exponent (restricted 3-vortex).}
\label{tab:power}
\begin{tabular}{@{}rc@{\hspace{2em}}rc@{}}
\toprule
$n$ & $\sigma_Q/|\bar{Q}|$ & $n$ & $\sigma_Q/|\bar{Q}|$ \\
\midrule
$-2.0$ & $5 \times 10^{-8}$  & $\mathbf{0.5}$ & $\mathbf{3 \times 10^{-11}}$ \\
$-1.0$ & $2 \times 10^{-9}$  & $1.0$ & $2 \times 10^{-10}$ \\
$-0.5$ & $8 \times 10^{-11}$ & $2.0$ & $1 \times 10^{-21}$ (exact) \\
$0.0$  & $5 \times 10^{-32}$ (trivial) & $3.0$ & $2 \times 10^{-8}$ \\
\bottomrule
\end{tabular}
\end{table}

In the restricted problem with $\Gamma_3 = \epsilon$, fractional variance scales as $\sigma_Q/|\bar{Q}| \propto \epsilon^{1.44}$, improving as the test vortex weakens.

\subsection{Exact member: $Q_{r^2} = \Gamma_\mathrm{tot} L_z - |P|^2$}

The identity follows algebraically:
\begin{equation}
Q_{r^2} = \sum_{i<j} \Gamma_i \Gamma_j\, r_{ij}^2 = \Gamma_\mathrm{tot} L_z - |P|^2,
\end{equation}
where $P = \sum \Gamma_i z_i$ is the conserved centre of vorticity. Since both $L_z$ and $|P|^2$ are conserved, $Q_{r^2}$ is exactly conserved. This is confirmed numerically with $R^2 = 0.9999999999$ across all configurations.

\subsection{Three-dimensional results}

For two coaxial vortex rings evolving under the Biot--Savart law, table~\ref{tab:3d} shows the conservation quality. $Q_{1/r}$ is an order of magnitude better conserved than~$Q_r$.

\begin{table}
\centering
\caption{Conservation of $Q_f$ for 3D vortex rings.}
\label{tab:3d}
\begin{tabular}{@{}lcc@{}}
\toprule
$f(r)$ & $\sigma_Q/|\bar{Q}|$ & Rank \\
\midrule
$\mathbf{1/r}$   & $\mathbf{3.78 \times 10^{-4}}$ & 1st \\
$\mathrm{e}^{-r}$ & $1.79 \times 10^{-3}$ & 2nd \\
$\sqrt{r}$        & $2.95 \times 10^{-3}$ & 3rd \\
$r$               & $4.36 \times 10^{-3}$ & 5th \\
\bottomrule
\end{tabular}
\end{table}

Under combined stretching and lateral motion, growth ratios confirm $Q_{1/r}$ as most robust: $Q_r$ grows by $41\times$ at stretch factor~4, while $Q_{1/r}$ grows by only $8.9\times$.

\subsection{The Green's function principle}

The optimal $Q_f$ in both dimensions uses the Green's function of the Laplacian: $G_2(r) = -\ln r / (2\pi)$ in 2D, $G_3(r) = 1/(4\pi r)$ in 3D. This is not coincidental. The kinetic energy of an inviscid incompressible flow can be written as \citep{saffman1992, majda2002}
\begin{equation}
E = \frac{1}{2} \int\!\!\int \boldsymbol{\omega}(\boldsymbol{x}) \cdot \boldsymbol{\omega}(\boldsymbol{y}) \, G_d(|\boldsymbol{x} - \boldsymbol{y}|) \,\mathrm{d}\boldsymbol{x}\,\mathrm{d}\boldsymbol{y}.
\label{eq:energy}
\end{equation}
For point vortices, this reduces to $Q_{G_d}$ up to self-interaction terms. Energy is exactly conserved by Euler's equations, so $Q_{G_d}$ inherits this property. I conjecture: \textit{in any spatial dimension~$d$, the optimal $Q_f$ uses $f(r) = G_d(r)$.}

\subsection{Optimal linear combination}

Using L-BFGS-B optimisation over 12 basis functions (powers, exponentials, logarithms), the best linear combination achieves fractional variance $5 \times 10^{-6}$ summed over all configurations---a $300$-fold improvement over the best single function ($\sqrt{r}$: total $1.59 \times 10^{-3}$). The optimal coefficients have alternating signs, indicating systematic cancellation.

\section{The $Q_f$ dichotomy}

\subsection{Concentration-detecting versus stretch-resistant}

For a Gaussian vortex blob with width~$\sigma$ and fixed circulation~$\Gamma$, the scaling $Q_f \propto \Gamma^2 \sigma^{2\alpha}$ reveals two regimes. Numerical measurements give:

\begin{table}
\centering
\caption{Concentration response exponent~$\alpha$ for different~$f$.}
\label{tab:dichotomy}
\begin{tabular}{@{}lcc@{}}
\toprule
$f(r)$ & $\alpha$ & Concentration response \\
\midrule
$-\ln r$    & $-0.63$ & Grows (detects) \\
$1/\sqrt{r}$ & $-0.47$ & Grows (detects) \\
$\sqrt{r}$  & $+0.49$ & Shrinks (blind) \\
$r$         & $+0.99$ & Shrinks (blind) \\
\bottomrule
\end{tabular}
\end{table}

Functions with $\alpha < 0$ (diverging~$f$ at small~$r$) grow as vorticity concentrates, providing a warning signal. If a concentration-detecting $Q_f$ is conserved, then concentration is bounded---directly constraining the Beale--Kato--Majda blowup criterion.

\subsection{Viscous decay rates}

Under Navier--Stokes evolution with viscosity~$\nu$, the members decay at different rates. $Q_{\sqrt{r}}$ has a linear decay relationship ($\nu^{0.99}$, $R^2 = 0.998$), giving a predictable lower bound
\begin{equation}
Q_{\sqrt{r}}(t) \geq Q_{\sqrt{r}}(0)\,\mathrm{e}^{-C\nu t}
\end{equation}
with measured constant $C \approx 7$ and coefficient of variation 5.6\%, compared to 46\% for $Q_{-\ln r}$. This makes $Q_{\sqrt{r}}$ the preferred choice for viscous regularity bounds.

\subsection{Curvature-weighted hybrid}

In three dimensions, curvature weighting $Q_{\kappa,f} = \int\!\!\int \kappa_i \kappa_j\, f(r_{ij})\,\mathrm{d}s\,\mathrm{d}t$ provides stretch resistance for any~$f$, because under stretching by factor~$s$: $\kappa \to \kappa/s$ and $\mathrm{d}s \to s\,\mathrm{d}s$, giving exact cancellation. Combining this with $f = 1/\sqrt{r}$ yields a quantity that is both stretch-resistant and concentration-detecting. Numerical tests on helical filaments confirm a factor-of-two improvement over $Q_{\kappa,r}$ (fractional variance 0.28 versus 0.60 at stretch factor~4).

\section{Negative result: higher-order invariants}

The full permutation sum $T = \sum_{i,j,k} \Gamma_i \Gamma_j \Gamma_k\,\mathrm{Area}(i,j,k)$ is identically zero by antisymmetry: the circulation product is symmetric under permutation while the signed area is antisymmetric, so each triplet contributes three positive and three negative terms that cancel. The ordered version $T_\mathrm{ord} = \sum_{i<j<k} \Gamma_i \Gamma_j \Gamma_k\,\mathrm{Area}(ijk)$ is not trivially zero but is not conserved: numerical tests give fractional variance of order $0.5$--$100$, compared to $< 10^{-7}$ for pairwise~$Q_f$. The pairwise structure is special and does not extend to higher orders.

\section{Discussion}

\subsection{Physical mechanism}

The time derivative is
\begin{equation}
\frac{\mathrm{d}Q_f}{\mathrm{d}t} = \sum_{i<j} \Gamma_i \Gamma_j\, f'(r_{ij})\,\frac{\mathrm{d}r_{ij}}{\mathrm{d}t}.
\end{equation}
Each $\mathrm{d}r_{ij}/\mathrm{d}t$ depends on the velocities of all vortices. But the velocity of vortex~$i$ is dominated by its strongest partner (the vortex with the largest $\Gamma_j / r_{ij}$), producing partial cancellation in the sum. The circulation weighting $\Gamma_i \Gamma_j$ reinforces this: strong pairs have nearly constant separation because they dominate each other's dynamics, while weak pair contributions are suppressed.

\subsection{Relation to known conservation laws}

The $Q_f$ family contains the Hamiltonian and angular impulse as exact members. The approximately conserved members interpolate between these and provide additional constraints beyond the four classical integrals. For numerical simulations, monitoring multiple $Q_f$ values provides a diagnostic of integrator quality beyond standard energy and momentum checks.

\subsection{Implications for Navier--Stokes regularity}

The dichotomy suggests a regularity strategy: (i)~choose $f = 1/\sqrt{r}$ (concentration-detecting); (ii)~prove $Q_f$ is bounded for smooth solutions; (iii)~boundedness implies vorticity concentration is limited, constraining the Beale--Kato--Majda criterion. The linear viscous decay of $Q_{\sqrt{r}}$ provides an additional ingredient: a predictable lower bound. Upgrading the numerical approximate conservation to rigorous bounds would be required for a full regularity argument; the present work provides candidate functionals.

\subsection{Limitations}

The $Q_f$ conservation is approximate, not exact (except for the two known members). Quality degrades with~$N$ (from $10^{-10}$ at $N\!=\!3$ to $10^{-4}$ at $N\!=\!8$) and depends on initial conditions. Rigorous error bounds on fractional variance as a function of~$N$, $T$, and the circulation distribution are not established here. The 3D results use discretised vortex filaments, introducing discretisation error.

\section{Conclusions}

A broad family of approximate conservation laws $Q_f = \sum_{i<j} \Gamma_i \Gamma_j\, f(r_{ij})$ has been identified for point vortex dynamics. The family contains two known exact invariants (energy and angular impulse) and infinitely many new approximate members. The optimal non-trivial member uses the Green's function of the Laplacian in any dimension---both giving the kinetic energy. The family exhibits a dichotomy between concentration-detecting and stretch-resistant members relevant to Navier--Stokes regularity. Higher-order (triplet) generalisations do not exist. These conservation laws provide new constraints for simulation validation, new diagnostic tools for vorticity concentration, and a framework for designing invariants tailored to specific physical questions.

\section*{Acknowledgements}

Numerical integration used SciPy's \texttt{solve\_ivp} (Dormand--Prince RK45). Optimisation used L-BFGS-B via SciPy. All computations performed on Apple Silicon (M3~Ultra). The author acknowledges the assistance of Claude (Anthropic) in developing the NoetherSolve framework, running numerical integrations, optimising invariants via L-BFGS-B, and assisting with manuscript preparation and \LaTeX\ formatting. All scientific content, derivations, interpretations, and final claims are the sole responsibility of the human author. Code and validation scripts: \texttt{https://github.com/SolomonB14D3/NoetherSolve}.

\section*{Declaration of interests}

The author reports no conflict of interest.

\section*{Data availability statement}

All data and code are available at \texttt{https://github.com/SolomonB14D3/NoetherSolve}.

\section*{Author ORCIDs}

B.~Sanchez, https://orcid.org/0009-0009-0421-1043.

\bibliographystyle{jfm}

\begin{thebibliography}{14}

\bibitem[\protect\citeauthoryear{Aref}{1983}]{aref1983}
{\sc Aref, H.} 1983
Integrable, chaotic, and turbulent vortex motion in two-dimensional flows.
{\it Annu.\ Rev.\ Fluid Mech.} {\bf 15}, 345--389.

\bibitem[\protect\citeauthoryear{Aref}{2007}]{aref2007}
{\sc Aref, H.} 2007
Point vortex dynamics: a classical mathematics playground.
{\it J.\ Math.\ Phys.} {\bf 48}~(6), 065401.

\bibitem[\protect\citeauthoryear{Dormand \& Prince}{1980}]{dormand1980}
{\sc Dormand, J.~R. \& Prince, P.~J.} 1980
A family of embedded {R}unge--{K}utta formulae.
{\it J.\ Comp.\ Appl.\ Math.} {\bf 6}~(1), 19--26.

\bibitem[\protect\citeauthoryear{Helmholtz}{1858}]{helmholtz1858}
{\sc Helmholtz, H.} 1858
{\"U}ber {I}ntegrale der hydrodynamischen {G}leichungen, welche den {W}irbelbewegungen entsprechen.
{\it J.\ reine angew.\ Math.} {\bf 55}, 25--55.

\bibitem[\protect\citeauthoryear{Kirchhoff}{1876}]{kirchhoff1876}
{\sc Kirchhoff, G.} 1876
{\it Vorlesungen {\"u}ber mathematische Physik: Mechanik}.
Teubner.

\bibitem[\protect\citeauthoryear{Majda \& Bertozzi}{2002}]{majda2002}
{\sc Majda, A.~J. \& Bertozzi, A.~L.} 2002
{\it Vorticity and Incompressible Flow}.
Cambridge University Press.

\bibitem[\protect\citeauthoryear{Modin \& Viviani}{2021}]{modin2021}
{\sc Modin, K. \& Viviani, M.} 2021
Integrability of point-vortex dynamics via symplectic reduction: a survey.
{\it Arnold Math.\ J.} {\bf 7}~(3), 357--385.

\bibitem[\protect\citeauthoryear{Newton}{2001}]{newton2001}
{\sc Newton, P.~K.} 2001
{\it The $N$-Vortex Problem: Analytical Techniques}.
Springer.

\bibitem[\protect\citeauthoryear{Newton}{2014}]{newton2014}
{\sc Newton, P.~K.} 2014
Point vortex dynamics in the post-{A}ref era.
{\it Fluid Dyn.\ Res.} {\bf 46}~(3), 031401.

\bibitem[\protect\citeauthoryear{Saffman}{1992}]{saffman1992}
{\sc Saffman, P.~G.} 1992
{\it Vortex Dynamics}.
Cambridge University Press.

\end{thebibliography}

\end{document}
